% ==== preamble.tex — pacchetti e macro comuni, senza \documentclass ====
% Incluso dal file master con % ==== preamble.tex — pacchetti e macro comuni, senza \documentclass ====
% Incluso dal file master con % ==== preamble.tex — pacchetti e macro comuni, senza \documentclass ====
% Incluso dal file master con % ==== preamble.tex — pacchetti e macro comuni, senza \documentclass ====
% Incluso dal file master con \input{preamble}

% --- Codifica e lingua (sillabazione italiana + lettere accentate) ---
% Nota: niente \usepackage[utf8]{inputenc}: dal 2018 l'UTF-8 è la codifica di
% default di LaTeX, e caricarlo a mano darebbe errore se un domani si passasse a
% xelatex/lualatex. Con pdflatex (il flusso di questo template) non serve.
\usepackage{lmodern} % Latin Modern: stesso disegno delle Computer Modern ma
                     % vettoriale -> caratteri più nitidi a schermo e a stampa
\usepackage[T1]{fontenc}
\usepackage[italian]{babel}

% =========================================================================
% LEGGIBILITÀ — le tre leve da toccare se il testo risulta faticoso.
% Sono indipendenti dal contenuto: cambiarle non rompe nulla, cambia solo
% l'impaginazione (e il numero di pagine).
%
%  1. CORPO DEL CARATTERE -> opzione del \documentclass in master.tex
%     (attualmente 11pt; 12pt per un testo ancora più grande).
%  2. LARGHEZZA DELLA RIGA -> i margini qui sotto. Righe troppo lunghe
%     stancano perché l'occhio fatica a trovare l'inizio della riga dopo:
%     l'ideale sta intorno ai 65-75 caratteri. Margini più larghi =
%     riga più corta. Con 11pt e 3cm si è intorno ai ~75.
%  3. INTERLINEA -> \setstretch più sotto (1 = default LaTeX, 1.15 =
%     circa "interlinea 1.15" dei word processor, 1.5 = molto arioso).
% =========================================================================

% --- Layout pagina: A4 ---
\usepackage[a4paper,
            top=2.5cm, bottom=2.5cm,
            left=3cm, right=3cm,
            includeheadfoot, headsep=0.6cm]{geometry}

% Interlinea (leva 3): 1.15 dà respiro senza sprecare pagine.
\usepackage{setspace}
\setstretch{1.15}

% Aria fra un paragrafo e il successivo: aiuta a distinguere i blocchi di
% testo a colpo d'occhio. Metti 0pt per tornare al comportamento LaTeX.
\setlength{\parskip}{0.35\baselineskip plus 0.1\baselineskip}

% Giustificazione più uniforme: riduce gli spazi larghi fra le parole e le
% sillabazioni, che è una causa nascosta di fatica nella lettura.
\usepackage{microtype}

\usepackage{graphicx}
\usepackage{amsmath}
\usepackage{amsfonts}
\usepackage{amssymb}
\usepackage{pgfplots}
\pgfplotsset{compat=1.18}
\usepackage{mathrsfs}
\usepackage{enumitem}
\usepackage{xcolor}
\usepackage{amsthm}
\usepackage{float}
\usepackage{needspace} % \Needspace*{Nh}: evita titoli isolati a fine pagina
\usepackage[most]{tcolorbox}
\usepackage[normalem]{ulem}
\usepackage{tikz}
\usetikzlibrary{positioning}
\usetikzlibrary{decorations.pathreplacing} % graffe: \draw[decorate, decoration={brace}]
\usepackage{fancyhdr}
% Nota: todonotes (con i marker \revisione e \figuramancante) va caricato DOPO
% hyperref — vedi più in basso, subito dopo \usepackage{hyperref}.

% =========================================================================
% FIGURE — ancoraggio \figura{<chiave>}{<descrizione>}
%
% Cerca figure/fig/<categoria>/<chiave>.tex fra le categorie elencate qui
% sotto: se lo trova lo include (centrato), altrimenti stampa il bollino blu
% \figuramancante con la descrizione.
%
% Conseguenza voluta: i file in lezioni/ contengono SOLO l'ancoraggio, mai il
% corpo della figura. Aggiungere un disegno = creare un file in figure/fig/,
% senza toccare la lezione; cancellare quel file = si torna da solo al
% bollino blu. Così ri-trascrivere una lezione non distrugge le figure, e
% generare figure non può corrompere il testo.
%
% La <chiave> è un identificatore stabile: NON deve dipendere dalla
% descrizione, che può cambiare a ogni ri-trascrizione.
%
% ATTENZIONE — dopo aver creato un NUOVO file figura serve `latexmk -g`
% (oppure due compilazioni): latexmk non può sapere che un file finora
% inesistente è diventato una dipendenza, e senza -g salta la ricompilazione
% lasciando il bollino blu anche se la figura c'è. Vale solo al primo
% passaggio da "assente" a "presente"; dopo, il file è tracciato normalmente.
%
% Per aggiungere una categoria basta aggiungere una riga #1{...} qui sotto.
% =========================================================================
\makeatletter
\newcommand{\figuracategorie}[1]{%
  #1{venn}%
  #1{grafi}%
  #1{grafici-funzioni}%
  #1{rette-numeriche}%
  #1{tabelle-verita}%
  #1{diagrammi-insiemistici}%
  #1{piano-cartesiano}%
  #1{altro}%
}
\newcommand{\fig@key}{}
\newcommand{\fig@path}{}
\newcommand{\fig@try}[1]{%
  \ifx\fig@path\@empty
    \edef\fig@cand{figure/fig/#1/\fig@key.tex}%
    \expandafter\IfFileExists\expandafter{\fig@cand}%
      {\global\let\fig@path\fig@cand}{}%
  \fi
}
\newcommand{\figura}[2]{%
  \gdef\fig@key{#1}%
  \gdef\fig@path{}%
  \figuracategorie{\fig@try}%
  \ifx\fig@path\@empty
    \figuramancante{#2}%
  \else
    \begin{center}\expandafter\input\expandafter{\fig@path}\end{center}%
  \fi
}
\makeatother
\usepackage{tocloft} % secondo indice: \newlistof per l'indice delle lezioni
% hyperref va caricato per ultimo (fra i pacchetti "normali").
% plainpages=false + pdfpagelabels: dà àncore di pagina univoche anche quando la
% numerazione riparte (frontespizio e primo foglio dell'indice), evitando il
% warning "duplicate destination page.*".
\usepackage[plainpages=false, pdfpagelabels]{hyperref}

% todonotes DOPO hyperref (ordine raccomandato dal pacchetto: così i marker e
% l'eventuale \listoftodos si agganciano correttamente ai link ipertestuali).
\usepackage[colorinlistoftodos, textsize=tiny]{todonotes} % marker \todo{REVIEW: ...}
% Marker rosso: contenuto di cui non si è sicuri della trascrizione, da verificare.
\newcommand{\revisione}[1]{\todo[inline, color=red!25]{REVIEW: #1}}
% Marker blu: nel materiale del professore c'è un disegno/schema da inserire a mano.
\newcommand{\figuramancante}[1]{\todo[inline, color=blue!25]{(inserire disegno) #1}}

\tcbuselibrary{theorems,breakable}

% Stile comune per tutti i blocchi
\tcbset{
    mystyle/.style={
        colback=gray!10,
        colframe=gray!50,
        fonttitle=\bfseries,
        coltitle=black,
        breakable,
        enhanced,
        before skip={\baselineskip},
        after skip={\baselineskip}
    }
}

% =========================================================================
\newtcbtheorem[number within=section]{theorem}{Teorema}{mystyle}{th}
\newtcbtheorem[number within=section]{property}{Proprietà}{mystyle}{pr}
\newtcbtheorem[number within=section]{definition}{Definizione}{mystyle}{def}
\newtcbtheorem[number within=section]{lemma}{Lemma}{mystyle}{lem}
% =========================================================================

% =========================================================================
% Secondo indice: "Indice delle lezioni" (oltre al normale \tableofcontents).
% Ogni lezione va marcata con \lezione{numero}{titolo} nel punto esatto in
% cui inizia (tipicamente subito prima del \section). La voce nell'indice
% rimanda con un collegamento ipertestuale a quel punto. Il titolo è
% facoltativo: \lezione{2}{} produce solo "Lezione 2".
\newlistof{lezioni}{loz}{Indice delle lezioni}
\newcommand{\lezione}[2]{%
    \phantomsection
    \addcontentsline{loz}{lezioni}{\textbf{Lezione~#1}\ifx\\#2\\\else\space--- #2\fi}%
}
\setlength{\cftlezioniindent}{0em}
\setlength{\cftlezioninumwidth}{0em}
\renewcommand{\cftlezionidotsep}{\cftdotsep} % puntini fino al numero di pagina
% =========================================================================

% --- Notazioni matematiche comuni (personalizza pure) ---
\newcommand{\R}{\mathbb{R}}
\newcommand{\N}{\mathbb{N}}
\newcommand{\Z}{\mathbb{Z}}
\newcommand{\Q}{\mathbb{Q}}
\newcommand{\C}{\mathbb{C}}

% --- Intestazione e piè di pagina ---
\pagestyle{fancy}
\fancyhf{}
\fancyhead[L]{\nouppercase{\leftmark}}
\fancyhead[R]{\thepage}
\renewcommand{\headrulewidth}{0.4pt}
\setlength{\headheight}{14pt}


% --- Codifica e lingua (sillabazione italiana + lettere accentate) ---
% Nota: niente \usepackage[utf8]{inputenc}: dal 2018 l'UTF-8 è la codifica di
% default di LaTeX, e caricarlo a mano darebbe errore se un domani si passasse a
% xelatex/lualatex. Con pdflatex (il flusso di questo template) non serve.
\usepackage{lmodern} % Latin Modern: stesso disegno delle Computer Modern ma
                     % vettoriale -> caratteri più nitidi a schermo e a stampa
\usepackage[T1]{fontenc}
\usepackage[italian]{babel}

% =========================================================================
% LEGGIBILITÀ — le tre leve da toccare se il testo risulta faticoso.
% Sono indipendenti dal contenuto: cambiarle non rompe nulla, cambia solo
% l'impaginazione (e il numero di pagine).
%
%  1. CORPO DEL CARATTERE -> opzione del \documentclass in master.tex
%     (attualmente 11pt; 12pt per un testo ancora più grande).
%  2. LARGHEZZA DELLA RIGA -> i margini qui sotto. Righe troppo lunghe
%     stancano perché l'occhio fatica a trovare l'inizio della riga dopo:
%     l'ideale sta intorno ai 65-75 caratteri. Margini più larghi =
%     riga più corta. Con 11pt e 3cm si è intorno ai ~75.
%  3. INTERLINEA -> \setstretch più sotto (1 = default LaTeX, 1.15 =
%     circa "interlinea 1.15" dei word processor, 1.5 = molto arioso).
% =========================================================================

% --- Layout pagina: A4 ---
\usepackage[a4paper,
            top=2.5cm, bottom=2.5cm,
            left=3cm, right=3cm,
            includeheadfoot, headsep=0.6cm]{geometry}

% Interlinea (leva 3): 1.15 dà respiro senza sprecare pagine.
\usepackage{setspace}
\setstretch{1.15}

% Aria fra un paragrafo e il successivo: aiuta a distinguere i blocchi di
% testo a colpo d'occhio. Metti 0pt per tornare al comportamento LaTeX.
\setlength{\parskip}{0.35\baselineskip plus 0.1\baselineskip}

% Giustificazione più uniforme: riduce gli spazi larghi fra le parole e le
% sillabazioni, che è una causa nascosta di fatica nella lettura.
\usepackage{microtype}

\usepackage{graphicx}
\usepackage{amsmath}
\usepackage{amsfonts}
\usepackage{amssymb}
\usepackage{pgfplots}
\pgfplotsset{compat=1.18}
\usepackage{mathrsfs}
\usepackage{enumitem}
\usepackage{xcolor}
\usepackage{amsthm}
\usepackage{float}
\usepackage{needspace} % \Needspace*{Nh}: evita titoli isolati a fine pagina
\usepackage[most]{tcolorbox}
\usepackage[normalem]{ulem}
\usepackage{tikz}
\usetikzlibrary{positioning}
\usetikzlibrary{decorations.pathreplacing} % graffe: \draw[decorate, decoration={brace}]
\usepackage{fancyhdr}
% Nota: todonotes (con i marker \revisione e \figuramancante) va caricato DOPO
% hyperref — vedi più in basso, subito dopo \usepackage{hyperref}.

% =========================================================================
% FIGURE — ancoraggio \figura{<chiave>}{<descrizione>}
%
% Cerca figure/fig/<categoria>/<chiave>.tex fra le categorie elencate qui
% sotto: se lo trova lo include (centrato), altrimenti stampa il bollino blu
% \figuramancante con la descrizione.
%
% Conseguenza voluta: i file in lezioni/ contengono SOLO l'ancoraggio, mai il
% corpo della figura. Aggiungere un disegno = creare un file in figure/fig/,
% senza toccare la lezione; cancellare quel file = si torna da solo al
% bollino blu. Così ri-trascrivere una lezione non distrugge le figure, e
% generare figure non può corrompere il testo.
%
% La <chiave> è un identificatore stabile: NON deve dipendere dalla
% descrizione, che può cambiare a ogni ri-trascrizione.
%
% ATTENZIONE — dopo aver creato un NUOVO file figura serve `latexmk -g`
% (oppure due compilazioni): latexmk non può sapere che un file finora
% inesistente è diventato una dipendenza, e senza -g salta la ricompilazione
% lasciando il bollino blu anche se la figura c'è. Vale solo al primo
% passaggio da "assente" a "presente"; dopo, il file è tracciato normalmente.
%
% Per aggiungere una categoria basta aggiungere una riga #1{...} qui sotto.
% =========================================================================
\makeatletter
\newcommand{\figuracategorie}[1]{%
  #1{venn}%
  #1{grafi}%
  #1{grafici-funzioni}%
  #1{rette-numeriche}%
  #1{tabelle-verita}%
  #1{diagrammi-insiemistici}%
  #1{piano-cartesiano}%
  #1{altro}%
}
\newcommand{\fig@key}{}
\newcommand{\fig@path}{}
\newcommand{\fig@try}[1]{%
  \ifx\fig@path\@empty
    \edef\fig@cand{figure/fig/#1/\fig@key.tex}%
    \expandafter\IfFileExists\expandafter{\fig@cand}%
      {\global\let\fig@path\fig@cand}{}%
  \fi
}
\newcommand{\figura}[2]{%
  \gdef\fig@key{#1}%
  \gdef\fig@path{}%
  \figuracategorie{\fig@try}%
  \ifx\fig@path\@empty
    \figuramancante{#2}%
  \else
    \begin{center}\expandafter\input\expandafter{\fig@path}\end{center}%
  \fi
}
\makeatother
\usepackage{tocloft} % secondo indice: \newlistof per l'indice delle lezioni
% hyperref va caricato per ultimo (fra i pacchetti "normali").
% plainpages=false + pdfpagelabels: dà àncore di pagina univoche anche quando la
% numerazione riparte (frontespizio e primo foglio dell'indice), evitando il
% warning "duplicate destination page.*".
\usepackage[plainpages=false, pdfpagelabels]{hyperref}

% todonotes DOPO hyperref (ordine raccomandato dal pacchetto: così i marker e
% l'eventuale \listoftodos si agganciano correttamente ai link ipertestuali).
\usepackage[colorinlistoftodos, textsize=tiny]{todonotes} % marker \todo{REVIEW: ...}
% Marker rosso: contenuto di cui non si è sicuri della trascrizione, da verificare.
\newcommand{\revisione}[1]{\todo[inline, color=red!25]{REVIEW: #1}}
% Marker blu: nel materiale del professore c'è un disegno/schema da inserire a mano.
\newcommand{\figuramancante}[1]{\todo[inline, color=blue!25]{(inserire disegno) #1}}

\tcbuselibrary{theorems,breakable}

% Stile comune per tutti i blocchi
\tcbset{
    mystyle/.style={
        colback=gray!10,
        colframe=gray!50,
        fonttitle=\bfseries,
        coltitle=black,
        breakable,
        enhanced,
        before skip={\baselineskip},
        after skip={\baselineskip}
    }
}

% =========================================================================
\newtcbtheorem[number within=section]{theorem}{Teorema}{mystyle}{th}
\newtcbtheorem[number within=section]{property}{Proprietà}{mystyle}{pr}
\newtcbtheorem[number within=section]{definition}{Definizione}{mystyle}{def}
\newtcbtheorem[number within=section]{lemma}{Lemma}{mystyle}{lem}
% =========================================================================

% =========================================================================
% Secondo indice: "Indice delle lezioni" (oltre al normale \tableofcontents).
% Ogni lezione va marcata con \lezione{numero}{titolo} nel punto esatto in
% cui inizia (tipicamente subito prima del \section). La voce nell'indice
% rimanda con un collegamento ipertestuale a quel punto. Il titolo è
% facoltativo: \lezione{2}{} produce solo "Lezione 2".
\newlistof{lezioni}{loz}{Indice delle lezioni}
\newcommand{\lezione}[2]{%
    \phantomsection
    \addcontentsline{loz}{lezioni}{\textbf{Lezione~#1}\ifx\\#2\\\else\space--- #2\fi}%
}
\setlength{\cftlezioniindent}{0em}
\setlength{\cftlezioninumwidth}{0em}
\renewcommand{\cftlezionidotsep}{\cftdotsep} % puntini fino al numero di pagina
% =========================================================================

% --- Notazioni matematiche comuni (personalizza pure) ---
\newcommand{\R}{\mathbb{R}}
\newcommand{\N}{\mathbb{N}}
\newcommand{\Z}{\mathbb{Z}}
\newcommand{\Q}{\mathbb{Q}}
\newcommand{\C}{\mathbb{C}}

% --- Intestazione e piè di pagina ---
\pagestyle{fancy}
\fancyhf{}
\fancyhead[L]{\nouppercase{\leftmark}}
\fancyhead[R]{\thepage}
\renewcommand{\headrulewidth}{0.4pt}
\setlength{\headheight}{14pt}


% --- Codifica e lingua (sillabazione italiana + lettere accentate) ---
% Nota: niente \usepackage[utf8]{inputenc}: dal 2018 l'UTF-8 è la codifica di
% default di LaTeX, e caricarlo a mano darebbe errore se un domani si passasse a
% xelatex/lualatex. Con pdflatex (il flusso di questo template) non serve.
\usepackage{lmodern} % Latin Modern: stesso disegno delle Computer Modern ma
                     % vettoriale -> caratteri più nitidi a schermo e a stampa
\usepackage[T1]{fontenc}
\usepackage[italian]{babel}

% =========================================================================
% LEGGIBILITÀ — le tre leve da toccare se il testo risulta faticoso.
% Sono indipendenti dal contenuto: cambiarle non rompe nulla, cambia solo
% l'impaginazione (e il numero di pagine).
%
%  1. CORPO DEL CARATTERE -> opzione del \documentclass in master.tex
%     (attualmente 11pt; 12pt per un testo ancora più grande).
%  2. LARGHEZZA DELLA RIGA -> i margini qui sotto. Righe troppo lunghe
%     stancano perché l'occhio fatica a trovare l'inizio della riga dopo:
%     l'ideale sta intorno ai 65-75 caratteri. Margini più larghi =
%     riga più corta. Con 11pt e 3cm si è intorno ai ~75.
%  3. INTERLINEA -> \setstretch più sotto (1 = default LaTeX, 1.15 =
%     circa "interlinea 1.15" dei word processor, 1.5 = molto arioso).
% =========================================================================

% --- Layout pagina: A4 ---
\usepackage[a4paper,
            top=2.5cm, bottom=2.5cm,
            left=3cm, right=3cm,
            includeheadfoot, headsep=0.6cm]{geometry}

% Interlinea (leva 3): 1.15 dà respiro senza sprecare pagine.
\usepackage{setspace}
\setstretch{1.15}

% Aria fra un paragrafo e il successivo: aiuta a distinguere i blocchi di
% testo a colpo d'occhio. Metti 0pt per tornare al comportamento LaTeX.
\setlength{\parskip}{0.35\baselineskip plus 0.1\baselineskip}

% Giustificazione più uniforme: riduce gli spazi larghi fra le parole e le
% sillabazioni, che è una causa nascosta di fatica nella lettura.
\usepackage{microtype}

\usepackage{graphicx}
\usepackage{amsmath}
\usepackage{amsfonts}
\usepackage{amssymb}
\usepackage{pgfplots}
\pgfplotsset{compat=1.18}
\usepackage{mathrsfs}
\usepackage{enumitem}
\usepackage{xcolor}
\usepackage{amsthm}
\usepackage{float}
\usepackage{needspace} % \Needspace*{Nh}: evita titoli isolati a fine pagina
\usepackage[most]{tcolorbox}
\usepackage[normalem]{ulem}
\usepackage{tikz}
\usetikzlibrary{positioning}
\usetikzlibrary{decorations.pathreplacing} % graffe: \draw[decorate, decoration={brace}]
\usepackage{fancyhdr}
% Nota: todonotes (con i marker \revisione e \figuramancante) va caricato DOPO
% hyperref — vedi più in basso, subito dopo \usepackage{hyperref}.

% =========================================================================
% FIGURE — ancoraggio \figura{<chiave>}{<descrizione>}
%
% Cerca figure/fig/<categoria>/<chiave>.tex fra le categorie elencate qui
% sotto: se lo trova lo include (centrato), altrimenti stampa il bollino blu
% \figuramancante con la descrizione.
%
% Conseguenza voluta: i file in lezioni/ contengono SOLO l'ancoraggio, mai il
% corpo della figura. Aggiungere un disegno = creare un file in figure/fig/,
% senza toccare la lezione; cancellare quel file = si torna da solo al
% bollino blu. Così ri-trascrivere una lezione non distrugge le figure, e
% generare figure non può corrompere il testo.
%
% La <chiave> è un identificatore stabile: NON deve dipendere dalla
% descrizione, che può cambiare a ogni ri-trascrizione.
%
% ATTENZIONE — dopo aver creato un NUOVO file figura serve `latexmk -g`
% (oppure due compilazioni): latexmk non può sapere che un file finora
% inesistente è diventato una dipendenza, e senza -g salta la ricompilazione
% lasciando il bollino blu anche se la figura c'è. Vale solo al primo
% passaggio da "assente" a "presente"; dopo, il file è tracciato normalmente.
%
% Per aggiungere una categoria basta aggiungere una riga #1{...} qui sotto.
% =========================================================================
\makeatletter
\newcommand{\figuracategorie}[1]{%
  #1{venn}%
  #1{grafi}%
  #1{grafici-funzioni}%
  #1{rette-numeriche}%
  #1{tabelle-verita}%
  #1{diagrammi-insiemistici}%
  #1{piano-cartesiano}%
  #1{altro}%
}
\newcommand{\fig@key}{}
\newcommand{\fig@path}{}
\newcommand{\fig@try}[1]{%
  \ifx\fig@path\@empty
    \edef\fig@cand{figure/fig/#1/\fig@key.tex}%
    \expandafter\IfFileExists\expandafter{\fig@cand}%
      {\global\let\fig@path\fig@cand}{}%
  \fi
}
\newcommand{\figura}[2]{%
  \gdef\fig@key{#1}%
  \gdef\fig@path{}%
  \figuracategorie{\fig@try}%
  \ifx\fig@path\@empty
    \figuramancante{#2}%
  \else
    \begin{center}\expandafter\input\expandafter{\fig@path}\end{center}%
  \fi
}
\makeatother
\usepackage{tocloft} % secondo indice: \newlistof per l'indice delle lezioni
% hyperref va caricato per ultimo (fra i pacchetti "normali").
% plainpages=false + pdfpagelabels: dà àncore di pagina univoche anche quando la
% numerazione riparte (frontespizio e primo foglio dell'indice), evitando il
% warning "duplicate destination page.*".
\usepackage[plainpages=false, pdfpagelabels]{hyperref}

% todonotes DOPO hyperref (ordine raccomandato dal pacchetto: così i marker e
% l'eventuale \listoftodos si agganciano correttamente ai link ipertestuali).
\usepackage[colorinlistoftodos, textsize=tiny]{todonotes} % marker \todo{REVIEW: ...}
% Marker rosso: contenuto di cui non si è sicuri della trascrizione, da verificare.
\newcommand{\revisione}[1]{\todo[inline, color=red!25]{REVIEW: #1}}
% Marker blu: nel materiale del professore c'è un disegno/schema da inserire a mano.
\newcommand{\figuramancante}[1]{\todo[inline, color=blue!25]{(inserire disegno) #1}}

\tcbuselibrary{theorems,breakable}

% Stile comune per tutti i blocchi
\tcbset{
    mystyle/.style={
        colback=gray!10,
        colframe=gray!50,
        fonttitle=\bfseries,
        coltitle=black,
        breakable,
        enhanced,
        before skip={\baselineskip},
        after skip={\baselineskip}
    }
}

% =========================================================================
\newtcbtheorem[number within=section]{theorem}{Teorema}{mystyle}{th}
\newtcbtheorem[number within=section]{property}{Proprietà}{mystyle}{pr}
\newtcbtheorem[number within=section]{definition}{Definizione}{mystyle}{def}
\newtcbtheorem[number within=section]{lemma}{Lemma}{mystyle}{lem}
% =========================================================================

% =========================================================================
% Secondo indice: "Indice delle lezioni" (oltre al normale \tableofcontents).
% Ogni lezione va marcata con \lezione{numero}{titolo} nel punto esatto in
% cui inizia (tipicamente subito prima del \section). La voce nell'indice
% rimanda con un collegamento ipertestuale a quel punto. Il titolo è
% facoltativo: \lezione{2}{} produce solo "Lezione 2".
\newlistof{lezioni}{loz}{Indice delle lezioni}
\newcommand{\lezione}[2]{%
    \phantomsection
    \addcontentsline{loz}{lezioni}{\textbf{Lezione~#1}\ifx\\#2\\\else\space--- #2\fi}%
}
\setlength{\cftlezioniindent}{0em}
\setlength{\cftlezioninumwidth}{0em}
\renewcommand{\cftlezionidotsep}{\cftdotsep} % puntini fino al numero di pagina
% =========================================================================

% --- Notazioni matematiche comuni (personalizza pure) ---
\newcommand{\R}{\mathbb{R}}
\newcommand{\N}{\mathbb{N}}
\newcommand{\Z}{\mathbb{Z}}
\newcommand{\Q}{\mathbb{Q}}
\newcommand{\C}{\mathbb{C}}

% --- Intestazione e piè di pagina ---
\pagestyle{fancy}
\fancyhf{}
\fancyhead[L]{\nouppercase{\leftmark}}
\fancyhead[R]{\thepage}
\renewcommand{\headrulewidth}{0.4pt}
\setlength{\headheight}{14pt}


% --- Codifica e lingua (sillabazione italiana + lettere accentate) ---
% Nota: niente \usepackage[utf8]{inputenc}: dal 2018 l'UTF-8 è la codifica di
% default di LaTeX, e caricarlo a mano darebbe errore se un domani si passasse a
% xelatex/lualatex. Con pdflatex (il flusso di questo template) non serve.
\usepackage{lmodern} % Latin Modern: stesso disegno delle Computer Modern ma
                     % vettoriale -> caratteri più nitidi a schermo e a stampa
\usepackage[T1]{fontenc}
\usepackage[italian]{babel}

% =========================================================================
% LEGGIBILITÀ — le tre leve da toccare se il testo risulta faticoso.
% Sono indipendenti dal contenuto: cambiarle non rompe nulla, cambia solo
% l'impaginazione (e il numero di pagine).
%
%  1. CORPO DEL CARATTERE -> opzione del \documentclass in master.tex
%     (attualmente 11pt; 12pt per un testo ancora più grande).
%  2. LARGHEZZA DELLA RIGA -> i margini qui sotto. Righe troppo lunghe
%     stancano perché l'occhio fatica a trovare l'inizio della riga dopo:
%     l'ideale sta intorno ai 65-75 caratteri. Margini più larghi =
%     riga più corta. Con 11pt e 3cm si è intorno ai ~75.
%  3. INTERLINEA -> \setstretch più sotto (1 = default LaTeX, 1.15 =
%     circa "interlinea 1.15" dei word processor, 1.5 = molto arioso).
% =========================================================================

% --- Layout pagina: A4 ---
\usepackage[a4paper,
            top=2.5cm, bottom=2.5cm,
            left=3cm, right=3cm,
            includeheadfoot, headsep=0.6cm]{geometry}

% Interlinea (leva 3): 1.15 dà respiro senza sprecare pagine.
\usepackage{setspace}
\setstretch{1.15}

% Aria fra un paragrafo e il successivo: aiuta a distinguere i blocchi di
% testo a colpo d'occhio. Metti 0pt per tornare al comportamento LaTeX.
\setlength{\parskip}{0.35\baselineskip plus 0.1\baselineskip}

% Giustificazione più uniforme: riduce gli spazi larghi fra le parole e le
% sillabazioni, che è una causa nascosta di fatica nella lettura.
\usepackage{microtype}

\usepackage{graphicx}
\usepackage{amsmath}
\usepackage{amsfonts}
\usepackage{amssymb}
\usepackage{pgfplots}
\pgfplotsset{compat=1.18}
\usepackage{mathrsfs}
\usepackage{enumitem}
\usepackage{xcolor}
\usepackage{amsthm}
\usepackage{float}
\usepackage{needspace} % \Needspace*{Nh}: evita titoli isolati a fine pagina
\usepackage[most]{tcolorbox}
\usepackage[normalem]{ulem}
\usepackage{tikz}
\usetikzlibrary{positioning}
\usetikzlibrary{decorations.pathreplacing} % graffe: \draw[decorate, decoration={brace}]
\usepackage{fancyhdr}
% Nota: todonotes (con i marker \revisione e \figuramancante) va caricato DOPO
% hyperref — vedi più in basso, subito dopo \usepackage{hyperref}.

% =========================================================================
% FIGURE — ancoraggio \figura{<chiave>}{<descrizione>}
%
% Cerca figure/fig/<categoria>/<chiave>.tex fra le categorie elencate qui
% sotto: se lo trova lo include (centrato), altrimenti stampa il bollino blu
% \figuramancante con la descrizione.
%
% Conseguenza voluta: i file in lezioni/ contengono SOLO l'ancoraggio, mai il
% corpo della figura. Aggiungere un disegno = creare un file in figure/fig/,
% senza toccare la lezione; cancellare quel file = si torna da solo al
% bollino blu. Così ri-trascrivere una lezione non distrugge le figure, e
% generare figure non può corrompere il testo.
%
% La <chiave> è un identificatore stabile: NON deve dipendere dalla
% descrizione, che può cambiare a ogni ri-trascrizione.
%
% ATTENZIONE — dopo aver creato un NUOVO file figura serve `latexmk -g`
% (oppure due compilazioni): latexmk non può sapere che un file finora
% inesistente è diventato una dipendenza, e senza -g salta la ricompilazione
% lasciando il bollino blu anche se la figura c'è. Vale solo al primo
% passaggio da "assente" a "presente"; dopo, il file è tracciato normalmente.
%
% Per aggiungere una categoria basta aggiungere una riga #1{...} qui sotto.
% =========================================================================
\makeatletter
\newcommand{\figuracategorie}[1]{%
  #1{venn}%
  #1{grafi}%
  #1{grafici-funzioni}%
  #1{rette-numeriche}%
  #1{tabelle-verita}%
  #1{diagrammi-insiemistici}%
  #1{piano-cartesiano}%
  #1{altro}%
}
\newcommand{\fig@key}{}
\newcommand{\fig@path}{}
\newcommand{\fig@try}[1]{%
  \ifx\fig@path\@empty
    \edef\fig@cand{figure/fig/#1/\fig@key.tex}%
    \expandafter\IfFileExists\expandafter{\fig@cand}%
      {\global\let\fig@path\fig@cand}{}%
  \fi
}
\newcommand{\figura}[2]{%
  \gdef\fig@key{#1}%
  \gdef\fig@path{}%
  \figuracategorie{\fig@try}%
  \ifx\fig@path\@empty
    \figuramancante{#2}%
  \else
    \begin{center}\expandafter\input\expandafter{\fig@path}\end{center}%
  \fi
}
\makeatother
\usepackage{tocloft} % secondo indice: \newlistof per l'indice delle lezioni
% hyperref va caricato per ultimo (fra i pacchetti "normali").
% plainpages=false + pdfpagelabels: dà àncore di pagina univoche anche quando la
% numerazione riparte (frontespizio e primo foglio dell'indice), evitando il
% warning "duplicate destination page.*".
\usepackage[plainpages=false, pdfpagelabels]{hyperref}

% todonotes DOPO hyperref (ordine raccomandato dal pacchetto: così i marker e
% l'eventuale \listoftodos si agganciano correttamente ai link ipertestuali).
\usepackage[colorinlistoftodos, textsize=tiny]{todonotes} % marker \todo{REVIEW: ...}
% Marker rosso: contenuto di cui non si è sicuri della trascrizione, da verificare.
\newcommand{\revisione}[1]{\todo[inline, color=red!25]{REVIEW: #1}}
% Marker blu: nel materiale del professore c'è un disegno/schema da inserire a mano.
\newcommand{\figuramancante}[1]{\todo[inline, color=blue!25]{(inserire disegno) #1}}

\tcbuselibrary{theorems,breakable}

% Stile comune per tutti i blocchi
\tcbset{
    mystyle/.style={
        colback=gray!10,
        colframe=gray!50,
        fonttitle=\bfseries,
        coltitle=black,
        breakable,
        enhanced,
        before skip={\baselineskip},
        after skip={\baselineskip}
    }
}

% =========================================================================
\newtcbtheorem[number within=section]{theorem}{Teorema}{mystyle}{th}
\newtcbtheorem[number within=section]{property}{Proprietà}{mystyle}{pr}
\newtcbtheorem[number within=section]{definition}{Definizione}{mystyle}{def}
\newtcbtheorem[number within=section]{lemma}{Lemma}{mystyle}{lem}
% =========================================================================

% =========================================================================
% Secondo indice: "Indice delle lezioni" (oltre al normale \tableofcontents).
% Ogni lezione va marcata con \lezione{numero}{titolo} nel punto esatto in
% cui inizia (tipicamente subito prima del \section). La voce nell'indice
% rimanda con un collegamento ipertestuale a quel punto. Il titolo è
% facoltativo: \lezione{2}{} produce solo "Lezione 2".
\newlistof{lezioni}{loz}{Indice delle lezioni}
\newcommand{\lezione}[2]{%
    \phantomsection
    \addcontentsline{loz}{lezioni}{\textbf{Lezione~#1}\ifx\\#2\\\else\space--- #2\fi}%
}
\setlength{\cftlezioniindent}{0em}
\setlength{\cftlezioninumwidth}{0em}
\renewcommand{\cftlezionidotsep}{\cftdotsep} % puntini fino al numero di pagina
% =========================================================================

% --- Notazioni matematiche comuni (personalizza pure) ---
\newcommand{\R}{\mathbb{R}}
\newcommand{\N}{\mathbb{N}}
\newcommand{\Z}{\mathbb{Z}}
\newcommand{\Q}{\mathbb{Q}}
\newcommand{\C}{\mathbb{C}}

% --- Intestazione e piè di pagina ---
\pagestyle{fancy}
\fancyhf{}
\fancyhead[L]{\nouppercase{\leftmark}}
\fancyhead[R]{\thepage}
\renewcommand{\headrulewidth}{0.4pt}
\setlength{\headheight}{14pt}
